%% ma: L12-B13-0023
%% lop: 12
%% chuong: 4
%% bai: 13
%% dang: 12.13.1
%% loai: TN4
%% muc: NB
%% dap_an: "D"
%% nguon: "(NBV)-ĐỀ SỐ 9-MỖI NGÀY 1 ĐỀ THI 2025.pdf (D09, MỖI NGÀY 1 ĐỀ - Đề số 9), Phần I Câu 11"
%% kien_thuc: "Bài 13 (thể tích khối tròn xoay khi quay hình phẳng quanh trục Ox)"
%% trang_thai: chinh-thuc
%% ghi_chu: "Đáp án gốc D đúng. Cùng ý tưởng với các câu dạng 12.13.1 (dạng đã duyệt)."
\begin{ex}
Cho hình phẳng $D$ giới hạn bởi đồ thị hàm số $y=f(x)=\cos\left(\dfrac x2+\dfrac\pi4\right)$, trục hoành và hai đường thẳng $x=0$ và $x=\pi$. Thể tích của khối tròn xoay tạo thành khi quay hình phẳng $D$ quanh trục hoành là
\choice{$V=\displaystyle\int_0^\pi\left|\cos\left(\dfrac x2+\dfrac\pi4\right)\right|\mathrm dx$}{$V=\displaystyle\int_0^\pi\cos^2\left(\dfrac x2+\dfrac\pi4\right)\mathrm dx$}{$V=\pi\displaystyle\int_0^\pi\left|\cos\left(\dfrac x2+\dfrac\pi4\right)\right|\mathrm dx$}{\True $V=\pi\displaystyle\int_0^\pi\cos^2\left(\dfrac x2+\dfrac\pi4\right)\mathrm dx$}
\loigiai{Thể tích khối tròn xoay sinh bởi hình phẳng giới hạn bởi đồ thị $y=f(x)$, trục hoành, $x=a$, $x=b$ khi quay quanh $Ox$ là $V=\pi\displaystyle\int_a^b f^2(x)\,\mathrm dx$. Với $f(x)=\cos\left(\dfrac x2+\dfrac\pi4\right)$, $a=0$, $b=\pi$ ta được phương án D.}
\end{ex}
